\section{Preliminaries and the Dyson equation}\label{sec:dyson}


We use the determinant and trace conventions introduced in Section~\ref{sec:introduction}.
In particular, all operator determinants use the normalized trace of the relevant matrix size.
We use multiplicativity and the block-diagonal determinant rule as in \cite[Proposition 2.3]{MS24}.

\subsection{The matrix Cauchy transform}
Consider first any selfadjoint pencil
\[
 K=A\otimes1+\sum_j B_j\otimes s_j\in M_d(\M),
 \qquad A=A^*,\quad B_j=B_j^*\in M_d(\C),
\]
Put $\Theta(b):=\sum_jB_jbB_j$ for $b\in M_d(\C)$.

\begin{definition}\label{def:cauchy-transform}
The matrix-valued Cauchy transform of $K$ is
\[
 G(w):=E_d[(w1_d-K)^{-1}],\quad \operatorname{Im}w>0.
\]
\end{definition}
We use the resolvent convention of \cite[Chapter 9, (9.5)]{MS17}.
The operator-valued semicircular moment rule gives the \emph{Dyson equation}
\begin{equation}\label{eq:dyson-general}
 G(w)^{-1}=w1_d-A-\Theta(G(w)),\qquad \operatorname{Im}G(w)<0.
\end{equation}
Here $\operatorname{Im}G:=(G-G^*)/(2i)$, and the inequality is negative definiteness.
The source input is \cite[Chapter 9, Definition 10, Theorem 11 and Proposition 13]{MS17}.
The moment recursion is
\[
 Q_0=1_d,\qquad Q_k=AQ_{k-1}
       +\sum_{r+s=k-2}\Theta(Q_r)Q_s,\quad Q_k:=E_d[K^k].
\]
The empty sum is zero.
Inserting this recursion into the convergent large-$|w|$ resolvent series proves \eqref{eq:dyson-general} there; analytic continuation proves it throughout the upper half-plane.

The imaginary-part condition selects exactly one solution at each $w$.
With $W:=iG$ and $V:=-i(w1_d-A)$, the equation takes the positive-real-part form $VW+\Theta(W)W=1_d$ studied in \cite[Theorem 2.1 and Section 3]{HRFS07} by holomorphic contraction.
We include a direct proof because its estimate also gives the invertibility needed when differentiating the Dyson equation.
Suppose $G_1,G_2$ are solutions and put $V_j:=-\operatorname{Im}G_j>0$, $v:=\operatorname{Im}w$.
Taking imaginary parts of the inverse equation gives
\begin{equation}\label{eq:imaginary-balance}
 V_j=G_j(v1_d+\Theta(V_j))G_j^*
     =G_j^*(v1_d+\Theta(V_j))G_j.
\end{equation}
For $D:=G_1-G_2$ we have $D=G_1\Theta(D)G_2$.
After setting
\[
 W:=V_1^{-1/2}DV_2^{-1/2},\quad
 L_i:=V_1^{-1/2}G_1B_iV_1^{1/2},\quad
 R_i:=V_2^{1/2}B_iG_2V_2^{-1/2},
\]
this becomes $W=\sum_iL_iWR_i$.
Equation~\eqref{eq:imaginary-balance} implies $\sum_iL_iL_i^*<1_d$ and $\sum_iR_i^*R_i<1_d$.
The row and column operator norm estimate therefore gives
\[
 \|W\|\leq
 \left\|\sum_iL_iL_i^*\right\|^{1/2}
 \left\|\sum_iR_i^*R_i\right\|^{1/2}\|W\|<\|W\|
\]
unless $W=0$.
This proves uniqueness.
We call this selected solution the \emph{resolvent solution}; no choice among real algebraic roots is implicit in the notation.
Taking $G_1=G_2=G$ in the same estimate shows that $X\mapsto G\Theta(X)G$ is a strict contraction in the norm $\|V^{-1/2}XV^{-1/2}\|$, where $V:=-\operatorname{Im}G$.
Consequently $X\mapsto X-G\Theta(X)G$ is invertible.
We use uniqueness to identify the variational critical points in Sections~\ref{sec:main-proof} and~\ref{sec:ordered-minmax}, and this invertibility for the derivatives in Section~\ref{subsec:brown-density}.

\subsection{Integrating the Cauchy transform}
The covariance is symmetric for the bilinear trace pairing: $\Tr(\Theta(X)Y)=\Tr(X\Theta(Y))$.
Choose the analytic logarithm of $\det G(w)$ compatible with $G(w)\sim w^{-1}1_d$ at infinity.
Then
\begin{equation}\label{eq:primitive}
 F(w):=-\frac1d\log\det G(w)
        +\frac1{2d}\Tr(\Theta(G(w))G(w))
      =\varphi_d(\Log(w1_d-K)).
\end{equation}
To prove this identity, differentiate \eqref{eq:dyson-general} and take the trace to obtain
\[
 \Tr G=-\Tr(G^{-1}G')+\Tr(\Theta(G)G').
\]
Thus both sides of \eqref{eq:primitive} have derivative $\tr_dG$.
Both are $\log w+o(1)$ as $w=it$ tends to infinity, which fixes the additive constant.
Taking real parts gives
\begin{equation}\label{eq:selfadjoint-potential}
 \log\Delta(K-w1_d)
 =-\frac1d\log|\det G(w)|
       +\frac1{2d}\operatorname{Re}\Tr(\Theta(G(w))G(w)).
\end{equation}

The identity extends to a real spectral gap by continuity.
It fixes the value at a resolvent solution; an extremum over free trial matrices requires a separate argument.
Related finite log-determinant expressions, evaluated at solutions of deterministic resolvent equations, occur in the random-matrix deterministic equivalents of \cite[Theorem 4.1]{HLN07}.
The shared method is integration of the resolvent; the minimization over arbitrary positive matrix trials is established in Section~\ref{sec:main-proof}.

\subsection{Selfadjoint blocks and regularization}\label{subsec:hermitization}
For the general pencil $S$ define its hermitization by
\begin{equation}\label{eq:hermitization}
 \mathscr H:=\begin{pmatrix}0&S\\S^*&0\end{pmatrix},\qquad
 A:=\begin{pmatrix}0&a_0\\a_0^*&0\end{pmatrix},\qquad
 B_j:=\begin{pmatrix}0&a_j\\a_j^*&0\end{pmatrix}.
\end{equation}
This is a selfadjoint pencil of size $2m$; see \cite[Section 3]{MS24} and \cite[Chapter 11, Section 9]{MS17} for this block construction.
Put $\Theta(b):=\sum_jB_jbB_j$ and
\[
 \eta^*(b):=\sum_j a_j^*ba_j,\qquad
 \rho^*(b):=\sum_j a_j^*ba_j^*.
\]
These are the Hilbert--Schmidt adjoints of $\eta$ and $\rho$ for $\langle b,c\rangle:=\Tr(b^*c)$.
We follow the notation of \cite[Lemma 3.5(ii)]{MS24}.
Direct multiplication gives
\begin{equation}\label{eq:full-covariance}
 \Theta\begin{pmatrix}b_{11}&b_{12}\\b_{21}&b_{22}\end{pmatrix}
 =\begin{pmatrix}\eta(b_{22})&\rho(b_{21})\\
                  \rho^*(b_{12})&\eta^*(b_{11})\end{pmatrix}.
\end{equation}
This displays the role of $\rho$ anticipated in the introduction.

For $\varepsilon>0$, set
\begin{equation}\label{eq:regularized-determinant}
 \log\Delta_\varepsilon(S)
       :=\tfrac12\varphi_m\log(S^*S+\varepsilon^21_m).
\end{equation}
The operators $S^*S$ and $SS^*$ have the same tracial spectral distribution, including at zero, by polar decomposition and traciality.
Since $\mathscr H^2=\operatorname{diag}(SS^*,S^*S)$, the resolvent $G_\varepsilon:=E_{2m}((i\varepsilon 1_{2m}-\mathscr H)^{-1})$ satisfies
\begin{equation}\label{eq:regularized-primitive}
 \log\Delta_\varepsilon(S)
 =-\frac1{2m}\log|\det G_\varepsilon|
  +\frac1{4m}\operatorname{Re}\Tr(\Theta(G_\varepsilon)G_\varepsilon).
\end{equation}
This follows from \eqref{eq:selfadjoint-potential}.
Monotone spectral integration gives $\log\Delta_\varepsilon(S)\downarrow\log\Delta(S)$, also when the limit is $-\infty$.

\subsection{Positive maps and spectral gaps}\label{subsec:gap-tools}
We record two tools for identifying spectral gaps from the matrix Cauchy transform.
For a linear map on a finite-dimensional space, write $r(\Psi)$ for its spectral radius, the largest modulus of its eigenvalues.
If $\Psi(b)=\sum_j c_jbc_j^*$ on $M_m(\C)$, then the row operator $[c_1\ \cdots\ c_n]$ gives
\[
 \|\Psi\|=\|\Psi(1_m)\|.
\]
Indeed it gives the upper bound, and evaluation at $1_m$ gives equality.
If $W>0$ and $W-\Psi(W)>0$, conjugating the map by $b\mapsto W^{1/2}bW^{1/2}$ makes its value at the identity strictly less than the identity.
Its spectral radius is therefore less than one.
In finite dimension this implies convergence of $(\id-\Psi)^{-1}=\sum_{k\ge0}\Psi^k$, which preserves positivity.
This is the finite-dimensional positive-map form of the strict Stein criterion; see \cite[Lemma 5.5 and Theorem 5.9]{Pop03} for its relation to similarity to a strict contraction.

Next let $K=K^*\in M_d(\M)$, with scalar spectral measure $\nu_K$, and let $g(w):=E_d((w1_d-K)^{-1})$.
For real $a<b$,
\begin{align}
 -\lim_{v\downarrow0}\frac1\pi
     \int_a^b\operatorname{Im}\tr_d g(x+iv)\dd x
   &=\nu_K((a,b))+\tfrac12\nu_K(\{a,b\}),
       \label{eq:stieltjes-interval}\\
 -\lim_{v\downarrow0}v\operatorname{Im}\tr_d g(x+iv)
   &=\nu_K(\{x\}).\label{eq:stieltjes-atom}
\end{align}
To verify the first formula, integrate the kernel $v/((x-t)^2+v^2)$ over $x\in(a,b)$ and then over $\nu_K$.
The normalized inner integral is bounded by one and tends to the indicator of $(a,b)$ with half weights at its endpoints.
The second formula follows in the same way from $v^2/((x-t)^2+v^2)$.
Dominated convergence proves both assertions.

Faithfulness implies $\supp\nu_K=\spec(K)$: every nonzero spectral projection has positive trace.
Consequently, if the matrix Cauchy transform $g$ continues analytically across a real interval and has selfadjoint values there, that interval is a spectral gap.
On its compact subintervals the imaginary part tends uniformly to zero, so \eqref{eq:stieltjes-interval} gives zero mass.
The converse follows from the operator resolvent.
This criterion concerns continuation of the matrix Cauchy transform, not an arbitrary solution of the matrix equation at a real parameter.
